\section{方法比较：简化流程不等于没有假设}
DPO 的吸引力在于它用一个更短的优化链条替代了 RLHF 的两阶段流程，但这不意味着它没有前提。它之所以成立，是因为我们接受了最优策略和 reward 之间的特定关系，并且愿意用 reference policy 与 KL 约束来保持训练稳定。

\begin{knowledgebox}{理解 DPO 时要保留的三点}
\begin{itemize}
  \item 它省掉的是显式 RM，不是偏好建模本身
  \item reference policy 仍然在目标里扮演稳定器角色
  \item 偏好数据质量依旧直接决定训练上限
\end{itemize}
\end{knowledgebox}

\subsection{本章小结}
DPO 的“优雅”来自推导简洁，但实验时仍然要关注数据质量、reference 选择和偏好噪声，不能把它理解成零成本替代品。

\newpage
